\chapter[Band 29: Kommutative Halbgruppen]{Band 29 auf einen Blick}
\noindent\textbf{Kommutative Halbgruppen}\par\medskip
Vorausgesetzt sind Halbgruppen, zweiseitige Teilbarkeit und
die elementaren Mengen- und Additionsgesetze. \(p\mid_A a\)
bezeichnet die zweiseitige Halbgruppenteilbarkeit.

\begin{BandUebersicht}
\textbf{Kommutativität}
\UebFormel{\begin{aligned}
&\CommutativeSemigroupAlg A\star\ \leftrightarrow\\
&\quad\SemigroupAlg A\star\land\\
&\quad\forall a,b\in A\ (a\star b=b\star a).
\end{aligned}}
\UebQuellen{\FormulaRefAuto{CommutativeSemigroupDef}[def]}
&
\textbf{Vertauschungsgesetze}
Für \(\CommutativeSemigroupAlg A\star\) und \(a,b,c,d\in A\):\UebFormel{\begin{aligned}
&(a\star b)\star c=(a\star c)\star b,\\
&(a\star b)\star(c\star d)\\
&\qquad=(a\star c)\star(b\star d).
\end{aligned}}
\UebQuellen{\FormulaRefAuto{CommutativeSemigroupThreeTermExchange}[thm];
\FormulaRefAuto{CommutativeSemigroupFourTermExchange}[thm]} \\
\addlinespace[.8ex]

\textbf{Addition und Vereinigung}
\UebFormel{\begin{aligned}
&+:\mathbb N^2\to\mathbb N,\\
&\cup:\powerset(M)^2\to\powerset(M).
\end{aligned}}
\UebQuellen{\FormulaRefAuto{NaturalAdditionSemigroup}[thm];
\FormulaRefAuto{SemigroupDef}[def]}
&
\textbf{Kommutative Beispiele}
\UebFormel{\begin{aligned}
&\CommutativeSemigroupAlg{\mathbb N}{+},\\
&\CommutativeSemigroupAlg{\powerset(M)}{\cup}.
\end{aligned}}
\UebQuellen{\FormulaRefAuto{NaturalAdditionCommutativeSemigroup}[thm];
\FormulaRefAuto{PowerSetUnionCommutativeSemigroup}[thm]} \\
\addlinespace[.8ex]

\textbf{Primelement ohne vorausgesetzte Einheit}
In einer kommutativen Halbgruppe ist \(p\) prim, wenn\UebFormel{\begin{aligned}
&p\in A,\quad\exists q\in A\ (p\nmid_A q),\\
&a,b\in A,\ p\mid_A(a\star b)\ \vdash\\
&\qquad p\mid_A a\lor p\mid_A b.
\end{aligned}}
\UebQuellen{\FormulaRefAuto{CommutativeSemigroupPrimeDef}[def]}
&
\textbf{Echtheit in Idealnotation}
Mit dem zweiseitigen Hauptideal \(J(p)\) gilt:\UebFormel{\begin{aligned}
&p,q\in A\ \vdash\\
&\quad(p\mid_A q\leftrightarrow q\in J(p)),\\
&p\text{ prim}\ \vdash\ J(p)\subsetneq A.
\end{aligned}}Die letzte Folgerung entfaltet die Echtheitsforderung.
\UebQuellen{\FormulaRefAuto{SemigroupTwoSidedDividesPrincipalIdeal}[thm];
\FormulaRefAuto{CommutativeSemigroupPrimeDef}[def]} \\
\addlinespace[.8ex]

\textbf{Homomorphismen und Isomorphismen}
Für zwei Strukturen derselben Klasse bedeutet ein Homomorphismus
\(H(f)\):\UebFormel{\begin{aligned}
&f:A\to B,\\
&x,y\in A\ \vdash\\
&\quad f(x\star y)=f(x)\diamond f(y).
\end{aligned}}Ein Isomorphismus \(I(f)\) ist ein bijektiver Homomorphismus:\UebFormel{\begin{aligned}
&I(f)\ \leftrightarrow H(f)\land f:A\bij B.
\end{aligned}}
\UebQuellen{\FormulaRefAuto{CommutativeSemigroupHomDef}[def];
\FormulaRefAuto{CommutativeSemigroupIsoDef}[def]}
&
\textbf{Komposition erhält die Struktur}
Für Homomorphismen \(f:A\to B\), \(g:B\to C\) zwischen
Strukturen dieser Klasse ist \(g\circ f\) wieder ein Homomorphismus.
Mit \(\bullet\) als Operation auf \(C\) gilt für \(x,y\in A\):\UebFormel{\begin{aligned}
&(g\circ f)(x\star y)\\
&\quad=(g\circ f)(x)\mathbin{\bullet}(g\circ f)(y).
\end{aligned}}Sind \(f,g\) Isomorphismen, so ist auch \(g\circ f\) ein Isomorphismus.
\UebQuellen{\FormulaRefAuto{SemigroupHomComposition}[thm];
\FormulaRefAuto{SemigroupIsoComposition}[thm]} \\
\addlinespace[.8ex]
\end{BandUebersicht}

